\chapter{System Architecture}\label{ch:architecture}
\thispagestyle{pagestyle}

\section{High-Level Overview}
\label{sec:arch-overview}

SHILATE is structured as three loosely coupled tiers, each running as an
independent process or container and communicating exclusively through a
Mosquitto MQTT broker. This design ensures that any tier can be replaced
or upgraded without modifying the others, and that the same communication
fabric used during training is available unchanged on the deployment
target.

\textbf{Tier 1 --- Unity Simulation.} A Unity 2022 LTS Editor instance
runs the vehicle physics, the obstacle course, and the sensor model. It
publishes observations and reward signals to MQTT and consumes control
commands from the same broker. The Unity process also hosts the
\textit{Training Controller} Editor window, which provides a
developer-facing UI for health monitoring and training control.

\textbf{Tier 2 --- Python RL Trainer.} A Python process wraps the MQTT
interface in a Gymnasium-compliant environment \cite{Towers2024},
instantiates a Stable-Baselines3 PPO agent \cite{Raffin2021}, and runs
the training loop. It emits structured telemetry back through MQTT so the
Editor window can display live metrics. After training, the same Python
process can be re-invoked in inference mode with a saved checkpoint.

\textbf{Tier 3 --- Eclipse Leda SDV Stack.} The Eclipse Leda SDV image
\cite{EclipseLeda2023} runs on QEMU or physical hardware. It hosts a
Mosquitto broker, a Kuksa databroker, an \texttt{mqtt-kuksa-feeder}
container that maps MQTT telemetry to VSS paths, and a Velocitas vehicle
application that monitors operational limits.

\section{Communication Backbone}
\label{sec:mqtt-topics}

All inter-tier communication passes through a single Mosquitto MQTT 3.1.1
broker \cite{OASIS2014}. The topic namespace uses the fixed prefix
\texttt{env0/}, reflecting the single-environment design. The complete
topic set is listed in \cref{table:mqtt-topics}.

\begin{table}[H]
\caption{MQTT topic map for the SHILATE training loop.}
\label{table:mqtt-topics}
\begin{tabular}{|p{8.5cm}|p{3cm}|p{2.5cm}|}
  \hline
  \textbf{Topic} & \textbf{Direction} & \textbf{Payload} \\
  \hline
  \texttt{env0/leda/control/steer}    & Python$\to$Unity & float $[-1,1]$ \\
  \hline
  \texttt{env0/leda/control/throttle} & Python$\to$Unity & float $[0,1]$ \\
  \hline
  \texttt{env0/leda/control/brake}    & Python$\to$Unity & float $[0,1]$ \\
  \hline
  \texttt{env0/leda/control/reset}    & Python$\to$Unity & ``1'' \\
  \hline
  \texttt{env0/vehicle/speed}         & Unity$\to$Python & float km/h \\
  \hline
  \texttt{env0/vehicle/training/obs}  & Unity$\to$Python & JSON array \\
  \hline
  \texttt{env0/vehicle/training/reward} & Unity$\to$Python & float \\
  \hline
  \texttt{env0/vehicle/training/done} & Unity$\to$Python & bool \\
  \hline
  \texttt{env0/vehicle/training/episode\_end} & Unity$\to$Editor & JSON object \\
  \hline
  \texttt{env0/vehicle/training/heartbeat}    & Unity$\to$Editor & JSON (1 Hz) \\
  \hline
\end{tabular}
\end{table}

QoS 0 (at-most-once delivery) is used throughout. Given that observations
arrive at 10~Hz and each new observation makes the previous one
irrelevant, occasional message loss is harmless and the absence of
acknowledgement handshakes reduces loop latency significantly compared to
QoS~1 or 2.

\section{Training Data-Flow}
\label{sec:dataflow}

Each training step follows a synchronised request-response pattern over
MQTT. The Python environment calls \texttt{step(action)}, which:
\begin{enumerate}[leftmargin=2cm,topsep=1.15pt,itemsep=1.15pt,partopsep=1.15pt,parsep=1.15pt]
  \item Publishes steer, throttle, and brake to the three control topics.
  \item Blocks on a threading event until Unity publishes a new
    \texttt{training/obs} message, or a 5-second timeout elapses.
  \item Reads the co-published reward and done values.
  \item Returns \texttt{(obs, reward, terminated, truncated, info)} to
    Stable-Baselines3.
\end{enumerate}
This cycle runs at the Unity physics rate (approximately 50~Hz
wall-clock), throttled by the 10~Hz sensor publication interval. The PPO
rollout buffer accumulates 2048 transitions before a gradient update.

\section{Configuration Management}
\label{sec:config}

Configuration is split across three locations:
\begin{itemize}[topsep=1.15pt,itemsep=1.15pt]
  \item \textbf{\texttt{config.json}} (repository root): holds
    \texttt{model.ray\_count} (default 9 in the root copy; the simulation
    runs 21 rays). Read by \texttt{config\_loader.py} with fallback.
  \item \textbf{\texttt{TrainingSettings.asset}} (Unity ScriptableObject):
    persists training knobs configured through the Editor window --- MQTT
    host/port, virtual-environment path, script paths, PPO
    hyperparameters, checkpoint save location, and optional resume-model
    path. Survives Editor restarts without an on-disk config file.
  \item \textbf{CLI flags} (\texttt{train.py}): hyperparameters from
    TrainingSettings are forwarded as command-line arguments when
    \texttt{PythonProcessManager} spawns the trainer, keeping the Python
    process stateless with respect to the Editor.
\end{itemize}

\section{Deployment Architecture on Eclipse Leda}
\label{sec:leda-arch}

When a trained checkpoint is deployed onto the Leda image, three
containers run simultaneously:
\begin{enumerate}[leftmargin=2cm,topsep=1.15pt,itemsep=1.15pt,partopsep=1.15pt,parsep=1.15pt]
  \item \textbf{\texttt{leda-controller}} (inference mode): loaded with
    the \texttt{best\_model.zip} checkpoint; subscribes to
    \texttt{env0/vehicle/training/obs} and publishes control commands.
  \item \textbf{\texttt{mqtt-kuksa-feeder}}: bridges
    \texttt{env0/vehicle/\#} telemetry topics to VSS paths in the Kuksa
    databroker via gRPC.
  \item \textbf{\texttt{velocitas-app}}: subscribes to
    \texttt{Vehicle.Speed} and engine RPM in the Kuksa databroker and
    publishes alerts to \texttt{leda/command/alert} when thresholds are
    exceeded.
\end{enumerate}
Each container is described by a Kanto manifest
(\texttt{kanto-manifest.json}) specifying the OCI image reference,
environment variables, and restart policy. The provisioning script
\texttt{leda/provision-device.sh} copies manifests to the device and
triggers Kanto to pull and start the containers.
